
\documentclass{article}
\usepackage{colortbl}
\usepackage{makecell}
\usepackage{multirow}
\usepackage{supertabular}

\begin{document}

\newcounter{utterance}

\centering \large Interaction Transcript for game `guesswhat', experiment `Abs\_Level\_3', episode 6 with reka{-}flash{-}3.1{-}t0.0.
\vspace{24pt}

{ \footnotesize  \setcounter{utterance}{1}
\setlength{\tabcolsep}{0pt}
\begin{supertabular}{c@{$\;$}|p{.15\linewidth}@{}p{.15\linewidth}p{.15\linewidth}p{.15\linewidth}p{.15\linewidth}p{.15\linewidth}}
    \# & $\;$A & \multicolumn{4}{c}{Game Master} & $\;\:$B\\
    \hline

    \theutterance \stepcounter{utterance}  
    & & \multicolumn{4}{p{0.6\linewidth}}{
        \cellcolor[rgb]{0.9,0.9,0.9}{
            \makecell[{{p{\linewidth}}}]{
                \texttt{\tiny{[P1$\langle$GM]}}
                \texttt{Let's play a guessing game. Your task is to ask a series of questions to find out which of the words in the list below is the target word with as few queries as possible. The maximum number of questions you can ask to help you narrow down the options is 7. Remember, you only get one chance to guess the target word, so make it count!} \\
\\ 
\texttt{The rules of the game are: } \\
\\ 
\texttt{1. Only ask one question per turn. DO NOT ASK MORE THAN ONE QUESTION.} \\
\\ 
\texttt{2. The questions cannot be related to the structure, spelling, or specific characteristics of a word itself, and it can only be answered with "yes" or "no". } \\
\\ 
\texttt{3. Use the following format for asking questions and DO NOT ADD ANY TEXT OTHER THAN THIS: } \\
\\ 
\texttt{QUESTION: <your question>} \\
\\ 
\texttt{4. You have only one opportunity to guess the target word. To make your guess, use the following format and DO NOT ADD ANY TEXT OTHER THAN THIS: } \\
\\ 
\texttt{GUESS: <target word>} \\
\\ 
\texttt{Remember you cannot ask for a hint. Make your question brief. Let’s start!} \\
\\ 
\texttt{List of candidate words: {[}'technician', 'optician', 'psychotherapist', 'chiropractor', 'surgeon', 'gp', 'nurse', 'pharmacist'{]}} \\
            }
        }
    }
    & \\ \\

    \theutterance \stepcounter{utterance}  
    & \multicolumn{4}{p{0.6\linewidth}}{
        \cellcolor[rgb]{0.9,0.9,0.9}{
            \makecell[{{p{\linewidth}}}]{
                \texttt{\tiny{[P1$\rangle$GM]}}
                \texttt{<reasoning>} \\
\texttt{Let's see. I need to figure out which word from the list is the target using up to 7 yes/no questions. The list is: {[}'technician', 'optician', 'psychotherapist', 'chiropractor','surgeon', 'gp', 'nurse', 'pharmacist'{]}.} \\
\\ 
\texttt{First, I should look for common characteristics among these words. All of them are medical professions, right? Let me check. Technician, optician, psychotherapist, chiropractor, surgeon, gp (general practitioner), nurse, pharmacist. Yep, all medical roles. So maybe the first question can narrow down based on that. But the rules say questions can't be about the structure or spelling, so I can't ask if it's a 9{-}letter word or something. So I need to think of other attributes.} \\
\\ 
\texttt{Since they're all medical, maybe the first question is about the type of role. For example, are they directly involved in patient care? Let's see: technician might be more support (like lab tech), optician makes glasses, psychotherapist is a mental health provider, chiropractor adjusts spine, surgeon operates, gp is primary care, nurse cares for patients, pharmacist dispenses meds. So some are direct care, others more support or administrative.} \\
\\ 
\texttt{Alternatively, maybe split into primary care vs specialists. GP is primary care. Surgeon is a specialist. But how to phrase that without mentioning medical terms? Wait, the question can't be about the structure or spelling, but can it be about their role? The rules say questions can't be related to the structure, spelling, or specific characteristics of the word itself. So maybe not about their job roles? Wait, the wording is a bit ambiguous. Let me check rule 2 again: "The questions cannot be related to the structure, spelling, or specific characteristics of a word itself, and it can only be answered with 'yes' or 'no'." So "specific characteristics of the word itself" probably refers to the word's form, like letters, syllables, parts of speech. So asking about their job roles is allowed, as that's not a characteristic of the word itself. So maybe that's okay.} \\
\\ 
\texttt{So first question: Is the target word a primary care provider? GP is primary care. Others like surgeon, chiropractor, etc., are specialists. So if the answer is yes, then GP is a candidate. If no,} \\
            }
        }
    }
    & & \\ \\

    \theutterance \stepcounter{utterance}  
    & & & \multicolumn{2}{p{0.3\linewidth}}{
        \cellcolor[rgb]{0.9,0.9,0.9}{
            \makecell[{{p{\linewidth}}}]{
                \texttt{\tiny{[GM$|$GM]}}
                \texttt{Invalid format. Guesser must use the form 'QUESTION: ' or 'GUESS: '.} \\
            }
        }
    }
    & & \\ \\

    \theutterance \stepcounter{utterance}  
    & & & \multicolumn{2}{p{0.3\linewidth}}{
        \cellcolor[rgb]{0.9,0.9,0.9}{
            \makecell[{{p{\linewidth}}}]{
                \texttt{\tiny{[GM$|$GM]}}
                \texttt{abort game} \\
            }
        }
    }
    & & \\ \\

\end{supertabular}
}

\end{document}
